%================================================================
\section{Results}\label{sec:results}

Every number in this section is a macro generated from the committed results in
\texttt{results/}. The section follows the order of the procedure: what history
matching kept, the image, the one-panel intervals at the three points, the price
of allowing for interference over repeated panels, and how the image moves with
the level tolerance.

\subsection{History matching kept a large set}\label{sec:nroy}

Against the planted targets, the three waves kept $\WaveKeptOne$, $\WaveKeptTwo$
and $\WaveKeptThree$ of their $\WavePoints$ points, and their union is
$|\NROY|=\NNroy$ (Figure~\ref{fig:waves} in Appendix~\ref{app:model}). The set
did not collapse to a point. The closest point has $I=\MinI$, far below
$\tau=\Tau$, yet it does not recover $\cstar$: its parameters are
$\alpha=\MinIAlpha$, $\beta=\MinIBeta$, $\gamma=\MinIGamma$, $\kappa=\MinIKappa$
and $\sigma_\varepsilon=\MinISigma$, against the planted values in
Table~\ref{tab:box}. A good fit does not pick out the truth. The median-I point has
$I=\MedI$ and a larger seeding effect, $\kappa=\MedIKappa$.

Which parameters do the targets restrict? For each parameter, the code divides
its standard deviation over the sampled NROY points by the standard deviation of a
uniform draw over the box on that axis. A ratio near one means the targets do not
restrict the parameter. The readout was declared before the run and is not a gate.

\begin{table}[t]
\centering
\caption{Spread ratios over the sampled NROY points.}
\label{tab:spread}
\begin{zebratabular}{lccccc}
\toprule
 & \colhead{$\alpha$} & \colhead{$\beta$} & \colhead{$\gamma$} & \colhead{$\kappa$} & \colhead{$\sigma_\varepsilon$} \\
\midrule
spread ratio & \SpreadRatioAlpha & \SpreadRatioBeta & \SpreadRatioGamma & \SpreadRatioKappa & \SpreadRatioSigma \\
\bottomrule
\end{zebratabular}
\tablenote{Standard deviation over the $\ThetaSample$ sampled NROY points divided
by the standard deviation of a uniform draw over $C$ on the same axis.}
\end{table}

Only the intercept $\alpha$ is pinned down (Table~\ref{tab:spread}). The ratios for
$\beta$, $\gamma$ and $\kappa$ are close to one, and $\sigma_\varepsilon$ is in
between. The pre-registration expected the peer-correlation moment to give a ratio
below one for $\gamma$. The ratio is $\SpreadRatioGamma$, below one only in the
third decimal, which is no restriction; that expectation was not met. It also
expected the pilot moment to restrict the trade-off between $\alpha$ and $\kappa$
more than the earlier design did. This paper does not rerun the earlier design,
so it cannot check that comparison; the sampled correlation between $\alpha$ and
$\kappa$ is $\SpreadCorrAlphaKappa$.

\subsection{The image of the policy effect is wide}\label{sec:image}

At the planted point, $\theta_1(\cstar)=\ThetaOneStar$ and
$\theta_2(\cstar)=\ThetaTwoStar$, so the premium of targeting is
$\prem(\cstar)=\PremiumStar$: in this village, seeding the hubs barely beats
seeding at random. That is a fact about this $\cstar$ and not a theorem about
targeting. The spillover is $\dspill(\cstar)=\SpillStar$.

Over the $\ThetaSample$ sampled NROY points, the image of the policy effect is
$\Theta_{\hm}=[\ThetaHmLo,\,\ThetaHmHi]$ and the image of the spillover is
$\Theta_{\expset}=[\ThetaExpLo,\,\ThetaExpHi]$. The report for $\theta_2$ that
the targets support is this range. It is neither a sampling distribution nor a
posterior. Figure~\ref{fig:points} shows every sampled point; for $\theta_2$ the
pointwise Monte Carlo error bars are smaller than the dots. A fresh recheck of
the four endpoint candidates moves the ends of $\Theta_{\hm}$ only in the sixth
decimal and the ends of $\Theta_{\expset}$ by amounts comparable to their Monte
Carlo errors, and it changes no verdict (Appendix~\ref{app:recheck}).

\begin{figure}[t]
\centering
\includefigure[width=0.92\linewidth]{image_points.png}
\caption{The $\ThetaSample$ sampled NROY points against their implausibility
$I(c)$. Left: $\theta_2(c)$. Right: $\dspill(c)$. Bars show $\pm1.96$ pointwise
Monte Carlo standard errors at each fixed $c$; dashed lines mark the image
endpoints.}
\label{fig:points}
\end{figure}

Why is the image wide? Largely because of the assumption $\delta$. With
$\delta_1=\DeltaOne$ on a level moment, the rule lets each level miss by up to
about $\sqrt{\tau\delta_1^2}$ on its own. That allowance exceeds
$\theta_1(\cstar)=\ThetaOneStar$, the very effect the pilot level carries, so the
pilot cannot pin the seeding effect $\kappa$. The spread ratio for $\kappa$,
$\SpreadRatioKappa$, is what that reading predicts, and the sensitivity menu of
Section~\ref{sec:menu} confirms it by tightening $\delta_1$. The peer weight is a
separate case: a check at $\cstar$ found that the peer-correlation moment hardly
moves across the box's range of $\gamma$; that check is not a committed result.

\subsection{One panel at each point}\label{sec:onepanel}

Table~\ref{tab:onepanel} reports every interval in the one panel drawn at each
diagnostic point, with both containment flags and the geometric verdict.
Figure~\ref{fig:overlay} draws the two foils against their images.

\begin{table}[t]
\centering
\caption{One panel at each diagnostic point. The first column is the field
experimenter; the other two are the calibrator. Each analyst row gives the
interval, then containment of the truth at the point (a), containment of the truth
at the village (b), and the verdict against the analyst's own image.}
\label{tab:onepanel}
\footnotesize
\setlength{\tabcolsep}{4pt}
\begin{zebratabular}{llll}
\toprule
 & \colhead{$\cstar$ (village)} & \colhead{min-I ($I=\MinI$)} & \colhead{median-I ($I=\MedI$)} \\
\midrule
\hiderowcolors\rowcolor{white}\multicolumn{4}{l}{\textit{Policy effect $\theta_2$; image $\Theta_{\hm}=[\ThetaHmLo,\,\ThetaHmHi]$}} \\
\rowcolor{white}truth $\theta_2(c_0)$ & $\ThetaTwoStar$ & $\MinIThetaTwo$ & $\MedIThetaTwo$ \\
\rowcolor{white}A0 (Foil 1) & $[\StarPtFoilOneLo,\,\StarPtFoilOneHi]$ & $[\MinIFoilOneLo,\,\MinIFoilOneHi]$ & $[\MedIFoilOneLo,\,\MedIFoilOneHi]$ \\
\rowcolor{white}\quad (a) / (b) / verdict & \StarPtFoilOneCoversTruth{} / \StarPtFoilOneCoversPlanted{} / \StarPtFoilOneVerdict & \MinIFoilOneCoversTruth{} / \MinIFoilOneCoversPlanted{} / \MinIFoilOneVerdict & \MedIFoilOneCoversTruth{} / \MedIFoilOneCoversPlanted{} / \MedIFoilOneVerdict \\
\rowcolor{zebra}A1 & $[\TradeoffStarAOneOnePanelLo,\,\TradeoffStarAOneOnePanelHi]$ & $[\TradeoffMinIAOneOnePanelLo,\,\TradeoffMinIAOneOnePanelHi]$ & $[\TradeoffMedIAOneOnePanelLo,\,\TradeoffMedIAOneOnePanelHi]$ \\
\rowcolor{zebra}\quad (a) / (b) / verdict & \TradeoffStarAOneOnePanelCoversTruth{} / \TradeoffStarAOneOnePanelCoversPlanted{} / \TradeoffStarAOneOnePanelVerdict & \TradeoffMinIAOneOnePanelCoversTruth{} / \TradeoffMinIAOneOnePanelCoversPlanted{} / \TradeoffMinIAOneOnePanelVerdict & \TradeoffMedIAOneOnePanelCoversTruth{} / \TradeoffMedIAOneOnePanelCoversPlanted{} / \TradeoffMedIAOneOnePanelVerdict \\
\rowcolor{white}A2 & $[\TradeoffStarATwoOnePanelLo,\,\TradeoffStarATwoOnePanelHi]$ & $[\TradeoffMinIATwoOnePanelLo,\,\TradeoffMinIATwoOnePanelHi]$ & $[\TradeoffMedIATwoOnePanelLo,\,\TradeoffMedIATwoOnePanelHi]$ \\
\rowcolor{white}\quad (a) / (b) / verdict & \TradeoffStarATwoOnePanelCoversTruth{} / \TradeoffStarATwoOnePanelCoversPlanted{} / \TradeoffStarATwoOnePanelVerdict & \TradeoffMinIATwoOnePanelCoversTruth{} / \TradeoffMinIATwoOnePanelCoversPlanted{} / \TradeoffMinIATwoOnePanelVerdict & \TradeoffMedIATwoOnePanelCoversTruth{} / \TradeoffMedIATwoOnePanelCoversPlanted{} / \TradeoffMedIATwoOnePanelVerdict \\
\midrule
\rowcolor{white}\multicolumn{4}{l}{\textit{Spillover $\dspill$; image $\Theta_{\expset}=[\ThetaExpLo,\,\ThetaExpHi]$}} \\
\rowcolor{white}truth $\dspill(c_0)$ & $\SpillStar$ & $\MinISpill$ & $\MedISpill$ \\
\rowcolor{white}Foil 2 & $[\StarPtFoilTwoLo,\,\StarPtFoilTwoHi]$ & $[\MinIFoilTwoLo,\,\MinIFoilTwoHi]$ & $[\MedIFoilTwoLo,\,\MedIFoilTwoHi]$ \\
\rowcolor{white}\quad (a) / (b) / verdict & \StarPtFoilTwoCoversTruth{} / \StarPtFoilTwoCoversPlanted{} / \StarPtFoilTwoVerdict & \MinIFoilTwoCoversTruth{} / \MinIFoilTwoCoversPlanted{} / \MinIFoilTwoVerdict & \MedIFoilTwoCoversTruth{} / \MedIFoilTwoCoversPlanted{} / \MedIFoilTwoVerdict \\
\bottomrule
\end{zebratabular}
\tablenote{Intervals at nominal level $\ConfidencePct\%$ from the one panel at
stream offset $0$. The flags describe that one panel; they are not coverage. At
$\cstar$ the verdict is reported for completeness and is not a test of the field
experimenter's estimator. Foil~2 targets $\dspill$ and is compared only with
$\Theta_{\expset}$.}
\end{table}

\begin{figure}[t]
\centering
\includefigure[width=0.86\linewidth]{overlay.png}
\caption{The two foils against their own images at the three diagnostic points
(rows). Left: $\theta_2$, the image $\Theta_{\hm}$ (band) and the A0 (Foil~1)
interval. Right: $\dspill$, the image $\Theta_{\expset}$ and the Foil~2 interval.
A star marks the truth at the point and a triangle the truth at $\cstar$ (omitted
in the $\cstar$ row). The two columns use different scales and different
functionals; they are never compared with each other.}
\label{fig:overlay}
\end{figure}

\paragraph{The field experimenter, at $\cstar$} Her A0 interval is
$[\StarPtFoilOneLo,\,\StarPtFoilOneHi]$ and the truth is $\ThetaTwoStar$, so the
interval misses; the point estimate is $\StarPtFoilOneBelowPct\%$ below the truth.
This failure needs no calibration. It is layer~(a), and it comes from the
identification conditions of the estimator: in this panel about
$\StarPtUnseededShareDisplayedPct\%$ of the gap is the effect of the campaign on
people who were not seeded, for which A0 has no term (Appendix~\ref{app:gap}). Her A1 and A2 intervals, computed from a pilot panel at
the same point, both contain $\theta_2(\cstar)$. Her Foil~2 interval contains the
spillover $\dspill(\cstar)$. Against the images, the A0, A1 and A2 intervals all
sit strictly inside $\Theta_{\hm}$ and the Foil~2 interval starts below the lower
end of $\Theta_{\expset}$, so it has no verdict. We report these verdicts and draw
no charge from them: she worked with the village's own panel.

\paragraph{The calibrator, at min-I and median-I} At both points the A0 interval
misses the truth at her own point (layer~a) and the truth at the village
(layer~b), and it lies strictly inside $\Theta_{\hm}$. Her A0 interval is
therefore wrong about her own world and also overconfident: it is far narrower than
the range of effects that the targets leave open. The A1 and A2 intervals behave
differently. At both points they contain the truth at her point. At the min-I
point they also contain the truth at the village, because the two truths are close
($\MinIThetaTwo$ and $\ThetaTwoStar$). At the median-I point they do not: the truth
there, $\MedIThetaTwo$, is far from the village's $\ThetaTwoStar$, and the
assumption ``my point is the world'' is costly. Even so, the A1 and A2 intervals
are F-OVERCONF at both points. They are right about her world and still narrower
than what the targets leave open. Foil~2 tells the same story for the spillover:
its interval contains both truths at both points and is F-OVERCONF against
$\Theta_{\expset}$ at both.

\paragraph{Containment and geometry are different tests} The A1, A2 and Foil~2
rows show why the two checks must be kept apart. An interval answers a question
about one world; the image answers a question about every world the targets cannot
rule out. A one-world interval can be right about its world and narrower than the
image at once. A calibrator who reports such an interval as her uncertainty about
the village understates what the targets leave open. At the min-I point she happens
to be right about the village as well; the median-I point shows that membership in
NROY does not guarantee it, and the image is what warns her.

\subsection{The price of allowing for interference}\label{sec:a2results}

Table~\ref{tab:tradeoff} reports the three analysts of $\theta_2$ over
$R=\TradeoffR$ fresh panels at each point, and Figure~\ref{fig:tradeoff} plots the
absolute bias against the mean interval width. No panel left A2 undefined: the
undefined count is $\TradeoffMinIATwoUndefinedCount$ at the min-I point,
$\TradeoffMedIATwoUndefinedCount$ at the median-I point and
$\TradeoffStarATwoUndefinedCount$ at $\cstar$.

\begin{table}[t]
\centering
\caption{The three analysts of $\theta_2$ over $\TradeoffR$ repeated panels at each
point (amendment A2).}
\label{tab:tradeoff}
\footnotesize
\setlength{\tabcolsep}{3.2pt}
\begin{zebratabular}{lcccccccc}
\toprule
 & \colhead{Bias} & \colhead{SD} & \colhead{RMSE} & \colhead{Mean SE} & \colhead{SE/SD}
 & \colhead{Containment [CP 95\%]} & \colhead{Width} & \colhead{Width/$\Theta_{\hm}$} \\
\midrule
\hiderowcolors\rowcolor{white}\multicolumn{9}{l}{\textit{min-I point, $\theta_2=\MinIThetaTwo$}} \\
\rowcolor{white}A0 & $\TradeoffMinIAZeroBias$ & $\TradeoffMinIAZeroSD$ & $\TradeoffMinIAZeroRMSE$ & $\TradeoffMinIAZeroMeanSE$ & $\TradeoffMinIAZeroSEOverSD$ & $\TradeoffMinIAZeroCoverage$ [$\TradeoffMinIAZeroCoverageLo$, $\TradeoffMinIAZeroCoverageHi$] & $\TradeoffMinIAZeroMeanWidth$ & $\TradeoffMinIAZeroWidthRatio$ \\
\rowcolor{zebra}A1 & $\TradeoffMinIAOneBias$ & $\TradeoffMinIAOneSD$ & $\TradeoffMinIAOneRMSE$ & $\TradeoffMinIAOneMeanSE$ & $\TradeoffMinIAOneSEOverSD$ & $\TradeoffMinIAOneCoverage$ [$\TradeoffMinIAOneCoverageLo$, $\TradeoffMinIAOneCoverageHi$] & $\TradeoffMinIAOneMeanWidth$ & $\TradeoffMinIAOneWidthRatio$ \\
\rowcolor{white}A2 & $\TradeoffMinIATwoBias$ & $\TradeoffMinIATwoSD$ & $\TradeoffMinIATwoRMSE$ & $\TradeoffMinIATwoMeanSE$ & $\TradeoffMinIATwoSEOverSD$ & $\TradeoffMinIATwoCoverage$ [$\TradeoffMinIATwoCoverageLo$, $\TradeoffMinIATwoCoverageHi$] & $\TradeoffMinIATwoMeanWidth$ & $\TradeoffMinIATwoWidthRatio$ \\
\midrule
\rowcolor{white}\multicolumn{9}{l}{\textit{median-I point, $\theta_2=\MedIThetaTwo$}} \\
\rowcolor{white}A0 & $\TradeoffMedIAZeroBias$ & $\TradeoffMedIAZeroSD$ & $\TradeoffMedIAZeroRMSE$ & $\TradeoffMedIAZeroMeanSE$ & $\TradeoffMedIAZeroSEOverSD$ & $\TradeoffMedIAZeroCoverage$ [$\TradeoffMedIAZeroCoverageLo$, $\TradeoffMedIAZeroCoverageHi$] & $\TradeoffMedIAZeroMeanWidth$ & $\TradeoffMedIAZeroWidthRatio$ \\
\rowcolor{zebra}A1 & $\TradeoffMedIAOneBias$ & $\TradeoffMedIAOneSD$ & $\TradeoffMedIAOneRMSE$ & $\TradeoffMedIAOneMeanSE$ & $\TradeoffMedIAOneSEOverSD$ & $\TradeoffMedIAOneCoverage$ [$\TradeoffMedIAOneCoverageLo$, $\TradeoffMedIAOneCoverageHi$] & $\TradeoffMedIAOneMeanWidth$ & $\TradeoffMedIAOneWidthRatio$ \\
\rowcolor{white}A2 & $\TradeoffMedIATwoBias$ & $\TradeoffMedIATwoSD$ & $\TradeoffMedIATwoRMSE$ & $\TradeoffMedIATwoMeanSE$ & $\TradeoffMedIATwoSEOverSD$ & $\TradeoffMedIATwoCoverage$ [$\TradeoffMedIATwoCoverageLo$, $\TradeoffMedIATwoCoverageHi$] & $\TradeoffMedIATwoMeanWidth$ & $\TradeoffMedIATwoWidthRatio$ \\
\midrule
\rowcolor{white}\multicolumn{9}{l}{\textit{$\cstar$ (village), $\theta_2=\ThetaTwoStar$}} \\
\rowcolor{white}A0 & $\TradeoffStarAZeroBias$ & $\TradeoffStarAZeroSD$ & $\TradeoffStarAZeroRMSE$ & $\TradeoffStarAZeroMeanSE$ & $\TradeoffStarAZeroSEOverSD$ & $\TradeoffStarAZeroCoverage$ [$\TradeoffStarAZeroCoverageLo$, $\TradeoffStarAZeroCoverageHi$] & $\TradeoffStarAZeroMeanWidth$ & $\TradeoffStarAZeroWidthRatio$ \\
\rowcolor{zebra}A1 & $\TradeoffStarAOneBias$ & $\TradeoffStarAOneSD$ & $\TradeoffStarAOneRMSE$ & $\TradeoffStarAOneMeanSE$ & $\TradeoffStarAOneSEOverSD$ & $\TradeoffStarAOneCoverage$ [$\TradeoffStarAOneCoverageLo$, $\TradeoffStarAOneCoverageHi$] & $\TradeoffStarAOneMeanWidth$ & $\TradeoffStarAOneWidthRatio$ \\
\rowcolor{white}A2 & $\TradeoffStarATwoBias$ & $\TradeoffStarATwoSD$ & $\TradeoffStarATwoRMSE$ & $\TradeoffStarATwoMeanSE$ & $\TradeoffStarATwoSEOverSD$ & $\TradeoffStarATwoCoverage$ [$\TradeoffStarATwoCoverageLo$, $\TradeoffStarATwoCoverageHi$] & $\TradeoffStarATwoMeanWidth$ & $\TradeoffStarATwoWidthRatio$ \\
\bottomrule
\end{zebratabular}
\tablenote{Bias is the mean estimate minus $\theta_2(c_0)$. SD is the standard
deviation of the estimates across panels (the sampling dispersion); Mean SE is the
mean reported standard error; Width is the mean width of the reported interval.
Containment is the share of panels whose interval contains $\theta_2(c_0)$, with an
exact Clopper--Pearson $95\%$ interval; for A1 and A2 it reflects conservative
intervals, not calibration.}
\end{table}

\begin{figure}[t]
\centering
\includefigure[width=\linewidth]{tradeoff.png}
\caption{Absolute bias against mean interval width over $\TradeoffR$ repeated
panels, one marker per analyst, one panel per diagnostic point. The dashed line
is the width of $\Theta_{\hm}$. The percentage beside each marker is the
repeated-sampling rate at which the analyst's interval contains $\theta_2(c_0)$,
printed as amendment A2 asked; for A1 and A2 it is conservative containment under
the Young bound and should not be read as calibrated coverage.}
\label{fig:tradeoff}
\end{figure}

\paragraph{Bias and spread} Absolute bias falls from A0 to A1 to A2 at every
point, and the sampling SD rises from A0 to A1 to A2 at every point. A0's bias is
large relative to its own spread: at $\cstar$ the bias is $\TradeoffStarAZeroBias$
with a sampling SD of $\TradeoffStarAZeroSD$, so its intervals almost never reach
the truth. A1 and A2 are biased downward at every point, by much less. The
amendment's proposed reason is two-hop spillover in the $\Sa$ control: people with
no seeded neighbor are still lifted by seeds two steps away, which raises the
control mean; the independent review of the code checked the sign of the bias and
not this mechanism. Allowing for interference does not remove the bias. It shrinks
it, and the analyst pays with dispersion. On mean squared error the trade does not
run all the way down the ladder: A1 has the smallest RMSE at every point, so A2's
further cut in bias costs more in spread than it saves.

\paragraph{The reported intervals are conservative} A0's reported standard error
matches its sampling SD closely (SE/SD of $\TradeoffMinIAZeroSEOverSD$,
$\TradeoffMedIAZeroSEOverSD$ and $\TradeoffStarAZeroSEOverSD$). For A1 and A2 the
reported SE is between $\TradeoffMinIATwoSEOverSD$ and $\TradeoffMedIAOneSEOverSD$
times the sampling SD. The A1 and A2 intervals are conservative relative to their
sampling distributions. The split of the reported variance that follows is
descriptive and is not a pre-registered result. It comes from addendum A2.1, which
was designed after the A2 results and the independent review of the code were
seen, states no expectation, and changes no estimator, interval or verdict
(Appendix~\ref{app:disclosures}). On the $\cstar$ panel, a share
\YoungShareAOneStar{} of A1's reported variance, and a share
\YoungShareATwoStar{} of A2's, is the Young bound on pairs with a zero
simulated joint exposure probability. The independent review of the code
(\texttt{reviews/a2-code.md}) traced the largest piece to the $\Sb$--$\Sa$ pair
rather than the $\Sone$--$\Sa$ pair. The $\pi>0$ term alone is smaller than the
design variance of the A1 estimate (fixed noise, \DesignRedrawsAOneStar{}
assignment redraws, mean main-term $\mathrm{SE}^2$ over design variance
\DesignMainTermRatioAOneStar; full mean $\mathrm{SE}^2$ over design variance
\DesignFullRatioAOneStar). The same pattern holds for both analysts at all three
points: the Young share runs from \YoungShareATwoMinI{} to \YoungShareAOneMedI,
the main-term ratio from \DesignMainTermRatioATwoMinI{} to
\DesignMainTermRatioAOneMinI, and the full ratio from \DesignFullRatioATwoMinI{}
to \DesignFullRatioAOneMedI. The Young terms are part of the pre-registered
variance and are not retuned here; a sharper variance would be a new design,
written before a new run.

\paragraph{Containment over repeated panels} A0's intervals contain
$\theta_2(c_0)$ in none of the panels at any point. A1's and A2's contain it in
every panel at every point, so the two are tied. That rate sits above the nominal
$\ConfidencePct\%$ and reflects the conservative variance; it is not evidence that
either interval is calibrated. Every repeated-sample mean width is narrower than
$\Theta_{\hm}$, the widest at $\TradeoffMedIATwoWidthRatio$ of its width. That is
interval geometry. It is not evidence that the interval is honest about what the
targets leave open, and it is not a coverage claim.

\paragraph{A post hoc reading of the trade-off} What follows is an
interpretation, formed after the A2 results and the A2.1 split were seen; it was
not pre-registered and it tests nothing. Moving from A0 to A1 and A2 does three
separate things. (i) It shrinks the bias but does not remove it. (ii) It raises the
sampling SD. This is a real rise in dispersion, because the SD is measured across
panels and does not depend on the variance estimator. (iii) It widens the reported interval by more
than the rise in SD alone would, because of the Young bound. Under the pilot's
design some pairs of people are never jointly exposed to the conditions that a
contrast compares, so the design cannot identify the covariance of their outcomes,
and the variance estimator replaces that covariance with Young's bound
\citep{aronow2017}. Part~(iii) depends on the exposure map. Foil~2's spillover
contrast uses the same pilot and the same simulated joint probabilities, yet on the
$\cstar$ panel only a share \YoungShareSpillStar{} of its variance is the bound
(Appendix~\ref{app:foil2}); A1 and A2 put weight on the $\Sb$ stratum, where most
of the bound's mass sits. The intervals that the bound widens are conservative,
and the containment rate of one comes from the bound. They are not calibrated.

\begin{table}[t]
\centering
\caption{The five expectations of amendment A2, written before any A2 number,
checked against Table~\ref{tab:tradeoff}.}
\label{tab:expect}
\begin{zebratabularx}{\linewidth}{lXl}
\toprule
\colhead{} & \colhead{Expectation (at every point)} & \colhead{Finding} \\
\midrule
(a) & $|$bias$|$ falls from A0 to A1 to A2 & met at all three points \\
(b) & sampling SD and width rise from A0 to A1 to A2 & met at all three points \\
(c) & A0 coverage far below nominal; A2 the closest to nominal & met only as a tie (A1 = A2) \\
(d) & A1 and A2 biased downward & met at all three points (sign only) \\
(e) & every interval narrower than $\Theta_{\hm}$ & met at all three points \\
\bottomrule
\end{zebratabularx}
\tablenote{(b): SD is the sampling dispersion and width is the conservative
interval; both rise, and they are different facts. (c): A0's rate is
$\TradeoffStarAZeroCoverage$ at every point; A1 and A2 both reach
$\TradeoffStarAOneCoverage$, so A2 is the closest only jointly with A1, and both
sit above nominal because their SE is between $\TradeoffMinIATwoSEOverSD$ and
$\TradeoffMedIAOneSEOverSD$ times the SD. (d): the review checked the sign; the
two-hop mechanism is the amendment's proposed reason. (e): mean width in
Table~\ref{tab:tradeoff}; the one-panel intervals are also strictly inside
$\Theta_{\hm}$ (Table~\ref{tab:onepanel}).}
\end{table}

Table~\ref{tab:expect} collects the five expectations. Four are met as written at
every point. Expectation~(c) is met only as a tie: A0's rate is far below nominal,
and A1 and A2 share the rate closest to nominal. The ladder therefore behaves as
the amendment predicted in its first-order features, bias down and spread up, and
it does not deliver a calibrated interval at any rung. More important for the
paper's question, no rung produces an interval as wide as the image. Even the
analyst who allows for interference and transports only within degree bins answers
a question about one world.

\subsection{How the image moves with the level tolerance}\label{sec:menu}

The image is a set under assumptions, and the assumption that appears to bind is
the level tolerance $\delta_1$. Amendment A1, written after the main results were
known (Appendix~\ref{app:disclosures}), reruns history matching and the image with
$\delta_1$ set to $\MenuDeltaMainLevel$ (the main design), $\MenuDeltaHalfLevel$ and
$\MenuDeltaQuarterLevel$; Appendix~\ref{app:menu} gives the tables. The kept set
shrinks from $\MenuDeltaMainNNroy$ points to $\MenuDeltaHalfNNroy$ and
$\MenuDeltaQuarterNNroy$, and at the smallest tolerance the image of the policy
effect narrows to $[\MenuDeltaQuarterThetaHmLo,\,\MenuDeltaQuarterThetaHmHi]$ while
the spread ratio for $\kappa$ falls to $\MenuDeltaQuarterSpreadRatioKappa$. That
supports the reading that the level allowance is what leaves $\kappa$ open. The
verdicts at $\cstar$ do not move at any entry. A smaller tolerance is a stronger
assumption, that the kernel matches the targets more closely; it buys a smaller set
and asks the reader to believe more. The menu is reported, not used to choose
$\delta$, and the results of this paper stay at the main design.
